\documentclass[11pt]{article}
\usepackage[utf8]{inputenc}
\usepackage[T1]{fontenc}
\usepackage{amsmath}
\usepackage{amssymb}
\usepackage{multicol}
\usepackage{geometry}
\usepackage{tikz}
\usepackage{enumitem}
\usepackage{xcolor}
\usepackage{titlesec}

% Configurações de layout
\geometry{a4paper, left=1cm, right=1cm, top=0.5cm, bottom=1.2cm}
\setlength{\columnseprule}{0.4pt}

% Cores para títulos
\titleformat{\section}{\normalfont\Large\bfseries\color{blue}}{\thesection}{1em}{}
\titleformat{\subsection}{\normalfont\large\bfseries\color{red}}{\thesubsection}{1em}{}

\title{\textcolor{red}{Primeiro Trimestre: Funções
    \\Função do Primeiro Grau - Atividade}}
\author{Professor: Jefferson}
\date{}

\begin{document}

\maketitle
\vspace{-1cm}

\begin{center}
\large{Nome: \underline{\hspace{8cm}} \quad Turma: \underline{\hspace{3cm}}}
\end{center}

\begin{multicols}{2}

\section*{Exercícios - Função do 1º Grau}

\begin{enumerate}[leftmargin=*]

\item Identifique o coeficiente angular ($a$) e o coeficiente linear ($b$) da função $f(x) = 3x + 7$.

\item Dada a função $f(x) = -2x + 4$, quais são os valores de $a$ e $b$?

\item Para a função $f(x) = 5x$, determine os coeficientes angular e linear.

\item Na função $f(x) = -x + 6$, identifique $a$ e $b$.

\item Calcule $f(2)$ para a função $f(x) = 4x - 3$.

\item Se $f(x) = 2x + 1$, qual é o valor de $f(3)$?

\item Para $f(x) = -3x + 5$, calcule $f(1)$.

\item Dada $f(x) = 6x + 2$, encontre $f(0)$.

\item Determine a raiz (zero) da função $f(x) = 2x - 8$.

\item Encontre a raiz da função $f(x) = -3x + 9$.

\item Qual é o zero da função $f(x) = 4x + 12$?

\item Para a função $f(x) = -5x + 10$, calcule a raiz.

\item Classifique a função $f(x) = 3x - 2$ em crescente ou decrescente.

\item A função $f(x) = -2x + 7$ é crescente ou decrescente? Justifique.

\item Determine se $f(x) = 5x + 1$ é crescente ou decrescente.

\item A função $f(x) = -x - 3$ é crescente ou decrescente?

\item Determine a função do primeiro grau que passa pelo ponto $(2, 5)$ e tem coeficiente angular $a = 3$.

\item Encontre a função que passa pelo ponto $(0, 4)$ com coeficiente angular $a = -1$.

\item Qual é a função que passa por $(1, 3)$ com $a = 2$?

\item Encontre a função que passa por $(3, 0)$ com coeficiente angular $a = -2$.

\item Determine a função que passa pelos pontos $(1, 2)$ e $(3, 8)$.

\item Encontre a função dados os pontos $(0, 5)$ e $(2, 1)$.

\item Qual é a função que passa por $(1, 1)$ e $(4, 4)$?

\item Determine a função que passa pelos pontos $(2, 3)$ e $(5, 9)$.

\item Para a função $f(x) = 2x - 4$, determine: a raiz, $f(0)$ e $f(3)$.

\item Faça uma análise completa de $f(x) = -3x + 6$: encontre a raiz, $f(0)$ e $f(1)$.

\item A temperatura em graus Celsius é dada por $T(t) = 20 + 5t$, onde $t$ é o tempo em horas. Qual é a temperatura inicial (em $t = 0$) e qual é a temperatura após 3 horas?

\item O custo de produção é dado por $C(x) = 50 + 10x$, onde $x$ é a quantidade produzida. Qual é o custo inicial e qual é o custo para produzir 50 unidades?

\item A altura de uma árvore é $h(t) = 1,5 + 0,3t$, onde $t$ é em anos. Qual é a altura inicial e qual é a altura após 10 anos?

\item Estude o sinal da função $f(x) = 2x - 4$, determinando onde $f(x) > 0$, $f(x) = 0$ e $f(x) < 0$.

\end{enumerate}

\end{multicols}

\end{document}
